\section{Introduction}
\label{sec:introduction}

Courts can force governments to deliver public services. They can
also change how the state produces those services. In right-to-health
litigation, a judge orders a public agency to provide a specific
medicine to an individual plaintiff. The agency must then source the
medicine quickly, often under short deadlines and one-sided penalties
for nondelivery. Court mandates do not merely change how much the
state pays; they change how the state is forced to buy.

This production margin is distinct from the usual fiscal question in
health litigation. A court order may have high private value and still
be costly to execute because it arrives outside the normal procurement
cycle. Conversely, a high observed price need not imply supplier
exploitation: it may reflect smaller lots, compressed search, or a
different supplier match. The empirical challenge is to distinguish
the legal mandate from the procurement mechanism through which the
mandate is fulfilled.

We study this question in S\~ao Paulo pharmaceutical procurement. The
key institutional fact is simple: court orders originate outside the
procurement process. Patients, physicians, lawyers, and judges
determine whether an order arrives, for which medicine, and when.
Purchasing units do not choose court-mandated status as a procurement
strategy. This institutional sequence supports interpreting urgent
purchases as externally imposed demand shocks to procurement,
conditional on item, time, and purchasing-unit fixed effects. This is
not a claim that patient or medicine selection into litigation is
ignorable. It is a claim about the procurement
transaction: the legal process moves demand into urgent execution from
outside the purchasing unit.

The identifying variation is institutional rather than experimental:
we compare procurement regimes for similar items, buyers, and periods
after a legal process has shifted demand from ordinary planning into
urgent execution.

The setting also offers a rare second comparison. Since
\BPsampleStartYear{}, S\~ao Paulo's health administration has operated
an administrative request channel through which some patients obtain
medicines outside court. These requests can generate urgent
procurement without judicial sanctions. The channel is useful because
it creates urgent pharmaceutical demand outside the court-sanction
regime. It is limited because requests are screened by a scientific
committee and are larger on average. We therefore treat
administrative-versus-litigated procurement as a bounded comparison
within urgent procurement, not as unselected sanction exposure.

The empirical strategy follows the institution.
Ordinary-versus-urgent comparisons estimate the procurement effect of
externally imposed legal urgency. Administrative-versus-litigated
comparisons ask how much of the within-urgent gap remains once
administrative selection is bounded using Lee trimming. Mechanism
tests then distinguish same-firm pricing from fragmented sourcing:
quantity is treated as a post-treatment channel, within
firm-buyer-item comparisons test whether incumbent suppliers charge
more for the same item to the same buyer, and winner-switching
evidence measures supplier-set reallocation.

The design is deliberately modest about what each comparison can
carry. The broad urgent comparison has the strongest institutional
claim. The administrative comparison is closer on the operational
environment but weaker on selection. The within-firm comparison is
sharp about pricing but silent about which suppliers win. Keeping
these claims separate avoids turning the administrative channel or the
mechanism regressions into more than they can support.

Four findings follow. First, urgent procurement is costlier and less
competitive: relative to ordinary purchases of comparable items by
the same buyers, urgent purchases have \BPnegPctHeadline{} higher
negotiated prices and \BPfirmsPctHeadlineAbs{} fewer bidders, while
tender success rises by \BPsuccessPP{}. Second, within urgent
procurement, court-mandated purchases are more expensive than
administrative urgent purchases. Because the administrative channel
is selected, the preferred estimate is a Lee-bounded
litigated-over-administrative gap of \BPutgBoundLow{} to
\BPutgBoundHigh{} rather than the naive gap of \BPutgPointNaive{}.
Third, there is little evidence of a broad same-firm markup in deep
repeated urgent markets: within firm-buyer-item triples observed
under both urgent regimes, the administrative coefficient is
\BPutgTripleCoefUTG{} (SE \BPutgTripleSEUTG{}). Fourth, the sourcing
margin is large. Administrative urgent orders are
\BPutgQtyAdminFoldPref{} times larger than litigated urgent orders;
among item-buyer pairs observed under both regimes, modal winners
differ in \BPwinnerSwitchDiffModalPct{} of pairs and mean winner-set
Jaccard similarity is \BPwinnerSwitchJaccardMean{}.

The contribution is to open the procurement black box of judicial
enforcement. Existing literatures study procurement waste,
bureaucratic accountability, courts, and right-to-health litigation.
This paper links them by separating two margins that standard
procurement comparisons confound: same-firm pricing and supplier-set
reallocation. In deep repeated urgent markets, court-mandated delivery
raises costs through fragmented sourcing rather than broad same-firm
markups. This reframes right-to-health litigation as a public-sector
production problem under legal constraint, connecting procurement
research on price dispersion, buyer capacity, and passive waste
\citep{bandiera2009active,best2023procurement,bosio2022public,
decarolis2018buyer}, courts and procurement performance
\citep{coviello2018court}, and accountability, bureaucratic
distortion, and health litigation
\citep{prendergast2007bureaucrats,bandiera2021allocation,
ferraz2009right,wang2015right,biehl2009will}. The policy implication
is not to weaken access to medicines. It is to preserve delivery while
rebuilding aggregation under legal urgency.
